\chapter{电阻与欧姆定律}\label{ch:g8:ohms-law}

电路的整个故事里，始终有一个问题贯穿其中：
有些元件让电流畅通无阻，有些则将电流限制成涓涓细流；
——\omterm{def:g3:first-electric-circuit:bulb}{灯泡}分得的\omterm{def:g8:voltage-voltmeter:voltage}{电压}不等，细导线发热，粗导线则保持凉爽。
阻碍本身如今有了名称和\omterm{def:g6:measurement-in-science:measuring}{单位}，
——在学完四年电路之后，迎来本书中最有用的定律：
三个字母，把推动力、电流与阻碍，
连成一行代数式。

\section{电阻}

\begin{definition}[电阻]\label{def:g8:ohms-law:resistance}
元件的\emph{电阻}\index{电阻} $R$ 衡量其
对电流的阻碍强弱：在\omterm{def:g8:voltage-voltmeter:voltage}{电压}一定时，电阻越大，
允许通过的电流就越弱。其\omterm{def:g6:measurement-in-science:measuring}{单位}为
\emph{欧姆}\index{欧姆}，
符号为 $\Omega$（希腊大写字母 omega），以纪念
这位中学教师兼物理学家，他以
自制导线和非凡耐心发现了本章定律。
\end{definition}

\begin{definition}[电阻器]\label{def:g8:ohms-law:resistor}
\emph{电阻器}\index{电阻器}是一种经制造、具有确定且稳定阻值的元件——
即经过标定的“障碍物”，
市售阻值从几分之一\omterm{def:g8:ohms-law:resistance}{欧姆}到数百万\omterm{def:g8:ohms-law:resistance}{欧姆}不等。电路利用\omterm{def:g8:ohms-law:resistance}{电阻}器
有意调节电流：压制会烧毁精密元件的电流，
用于调光、分压与调节。
其色环直接标明阻值；其符号
是一个简单的矩形。
\end{definition}

\begin{example}[身边的电阻]\label{ex:g8:ohms-law:values}
一米粗铜线：百分之几\omterm{def:g8:ohms-law:resistance}{欧姆}——平坦的大路。
发光的手电筒\omterm{def:g3:first-electric-circuit:bulb}{灯泡}：几\omterm{def:g8:ohms-law:resistance}{欧姆}。小型
\omterm{ex:g6:circuits-and-safety:motor}{电动机}：几\omterm{def:g8:ohms-law:resistance}{欧姆}。电热水壶的加热丝：约
$25\,\Omega$。干燥的人体皮肤，两手之间：数万
\omterm{def:g8:ohms-law:resistance}{欧姆}——湿时骤降：古老的“湿手”规则
正是一条关于\omterm{def:g8:ohms-law:resistance}{电阻}的陈述。断开的\omterm{ex:g3:first-electric-circuit:switch}{开关}：\omterm{def:g8:ohms-law:resistance}{电阻}
大到无法\omterm{def:g6:measurement-in-science:measuring}{测量}——无穷大的障碍。
\end{example}

\section{实验}

\begin{method}[测量元件的 $U$--$I$ 特性曲线]\label{met:g8:ohms-law:portrait}
一个\omterm{def:g8:ohms-law:resistor}{电阻器}，一台可调电源（或一串逐节增加的
\omterm{def:g3:first-electric-circuit:battery}{电池}），两只电表：
\begin{enumerate}
  \item 将电源、\omterm{def:g8:intensity-ammeter:ammeter}{电流表}和\omterm{def:g8:ohms-law:resistor}{电阻器}\omterm{def:g4:series-circuits:series}{串联}，
        \omterm{def:g8:voltage-voltmeter:voltmeter}{电压表}\omterm{def:g5:series-parallel-circuits:parallel}{并联}在\omterm{def:g8:ohms-law:resistor}{电阻器}两端；
  \item 将\omterm{def:g8:voltage-voltmeter:voltage}{电压}调到最小，记录数据对（$U$，$I$）；
  \item 逐级升高\omterm{def:g8:voltage-voltmeter:voltage}{电压}——一节\omterm{def:g3:first-electric-circuit:battery}{电池}、两节、
        三节，每次各记一组数据；
  \item 描点，以 $I$ 为横轴、$U$ 为纵轴，然后观察。
\end{enumerate}
\end{method}

\begin{example}[电阻器的“自白”]\label{ex:g8:ohms-law:confession}
某次实验的数据表，针对同一\omterm{def:g8:ohms-law:resistor}{电阻器}：
\begin{center}
\begin{tabular}{l|cccc}
$U$ (\unit{V}) & $1.5$ & $3.0$ & $4.5$ & $6.0$ \\
\hline
$I$ (\unit{A}) & $0.10$ & $0.20$ & $0.30$ & $0.40$
\end{tabular}
\end{center}
\omterm{def:g8:voltage-voltmeter:voltage}{电压}加倍，电流也加倍；\omterm{def:g8:voltage-voltmeter:voltage}{电压}增至三倍，电流也增至三倍：$U$ 与 $I$
成\emph{正比}，每个比值 $U/I$ 都
恒为同一个数：$15$。这一恒定比值并非偶然，它正是
\omterm{def:g8:ohms-law:resistor}{电阻器}的\omterm{def:g8:ohms-law:resistance}{电阻} $R = 15\,\Omega$，
在每一行数据中都显现出来。
\end{example}

\begin{proposition}[欧姆定律]\label{prop:g8:ohms-law:law}
对于\omterm{def:g3:temperature-thermometers:temperature}{温度}保持不变的\omterm{def:g8:ohms-law:resistor}{电阻器}，其两端的\omterm{def:g8:voltage-voltmeter:voltage}{电压}
与通过它的电流成正比，比例常数
即为其\omterm{def:g8:ohms-law:resistance}{电阻}：
\[
U = R \times I ,
\]
其中 $U$ 以\omterm{def:g8:voltage-voltmeter:voltage}{伏特}为\omterm{def:g6:measurement-in-science:measuring}{单位}，$I$ 以\omterm{def:g8:intensity-ammeter:intensity}{安培}为\omterm{def:g6:measurement-in-science:measuring}{单位}，$R$ 以\omterm{def:g8:ohms-law:resistance}{欧姆}为\omterm{def:g6:measurement-in-science:measuring}{单位}。变形后
得到：$I = U/R$（\omterm{def:g8:voltage-voltmeter:voltage}{电压}推动电流通过障碍的大小）和
$R = U/I$（由一对可靠读数所揭示的
障碍之值）。
\end{proposition}

\begin{omfigure}
\begin{tikzpicture}
\begin{axis}[omaxis, width=10.2cm, height=6.0cm,
    xlabel={$I$ (\unit{A})}, ylabel={$U$ (\unit{V})},
    xtick={0,0.1,0.2,0.3,0.4}, ytick={0,1.5,3.0,4.5,6.0},
    ymin=0, ymax=6.8, xmin=0, xmax=0.46]
  \addplot[very thick, omDef] coordinates {(0,0) (0.44,6.6)};
  \addplot[only marks, mark=*, omThm] coordinates
    {(0.10,1.5) (0.20,3.0) (0.30,4.5) (0.40,6.0)};
  \node[right, font=\small, omDef] at (axis cs:0.06,3.4)
    {过原点的直线};
\end{axis}
\end{tikzpicture}

{\small \omterm{def:g8:ohms-law:resistor}{电阻器}的 $U$--$I$ 特性曲线：各\omterm{def:g6:measurement-in-science:measuring}{测量}点落在
一条过原点的直线上——正比关系已画出。
直线越陡，\omterm{def:g8:ohms-law:resistance}{电阻}越大。}
\end{omfigure}

\begin{example}[定律可算出一切]\label{ex:g8:ohms-law:compute}
已知 $R$ 后，定律可从各个方向作答。
电热水壶的加热丝，$25\,\Omega$，接在 \qty{230}{V} 市电上：
$I = U/R = 230 \div 25 = \qty{9.2}{A}$——十\omterm{def:g8:intensity-ammeter:intensity}{安培}用电器
的来历就此解释清楚。一个\omterm{def:g8:ohms-law:resistor}{电阻器}须将精密\omterm{def:g3:first-electric-circuit:bulb}{灯泡}的
电流限制在 \qty{0.02}{A}，电源\omterm{def:g8:voltage-voltmeter:voltage}{电压}为 \qty{4.5}{V}：$R = U/I = 4.5
\div 0.02 = 225\,\Omega$——选取下一个标准阻值。
通过 $60\,\Omega$ 的电流为 \qty{0.05}{A}，其两端\omterm{def:g8:voltage-voltmeter:voltage}{电压}
为 $U = 60 \times 0.05 = \qty{3}{V}$。一条定律，三种
算法，足以应付全部计算。
\end{example}

\begin{remark}[三角辅助记法]\label{rem:g8:ohms-law:triangle}
有些学习者在一个小三角形中用铅笔标出 $U$、$R$ 和 $I$，
遮住所求字母，即可读出公式。
作为拐杖，这并无害处——但你的代数运算现在可以
在一行之内诚实地变形 $U = R I$，而这种诚实的做法适用于
你将遇到的每一个公式，无论是否有三角形。今年若需要，
不妨依靠这根拐杖；但要计划着学会独立行走。
\end{remark}

\section{电阻解释了什么}

\begin{example}[旧谜题，一份账单全解决]\label{ex:g8:ohms-law:mysteries}
定律一章中\omterm{def:g8:voltage-voltmeter:voltage}{电压}的不均分配：\omterm{def:g4:series-circuits:series}{串联}\omterm{def:g3:first-electric-circuit:bulb}{灯泡}
共用同一电流 $I$，由 $U = R I$ 可知，\omterm{def:g8:ohms-law:resistance}{电阻} $R$ 越大，
分得的\omterm{def:g8:voltage-voltmeter:voltage}{电压} $U$ 越大——出力更多者的较大\omterm{def:g8:voltage-voltmeter:voltage}{电压}降，如今可以算出。
\omterm{def:g5:series-parallel-circuits:parallel}{并联}一章中被节流的\omterm{def:g5:series-parallel-circuits:parallel}{支路}：那根细长导线
增加了\omterm{def:g8:ohms-law:resistance}{电阻}，$I = U/R$ 使该\omterm{def:g5:series-parallel-circuits:parallel}{支路}的分流减小。
\omterm{def:g7:short-circuits-safety:device}{短路}：旁路的 $R$ 几乎为零，于是 $I = U/R$
猛增——那场“暴冲”原来一直是除以近似于零的数。
四年的定性电路知识，一条定律的账单全部结算。
\end{example}

\begin{example}[导线的电阻，以米计]\label{ex:g8:ohms-law:wires}
导线的\omterm{def:g8:ohms-law:resistance}{电阻}随长度增加而增大，随粗细增加而减小——
细长导线好比崎岖的山路，
短粗导线则好比高速公路。因此，电工为耗电大的用电器配粗电缆，
加热线圈则特意采用又长又细的
\omterm{def:g8:ohms-law:resistance}{电阻}合金——同一条定律既使\omterm{def:g3:first-electric-circuit:bulb}{灯泡}发光，也使烤面包机发热，
其缘由在于：\omterm{def:g8:ohms-law:resistance}{电阻}是推动力耗费之处，
而耗费掉的推动力，正如\omterm{def:g5:everyday-energy:energy}{能量}一章明年将证实的那样，
化为\omterm{def:g5:heat-and-insulation:heat}{热量}。
\end{example}

\begin{remark}[定律的适用范围]\label{rem:g8:ohms-law:limits}
\omterm{def:g8:ohms-law:resistance}{欧姆}定律是\omterm{def:g8:ohms-law:resistor}{电阻器}的定律，并非普遍成立的规律。
\omterm{def:g6:measurement-in-science:measuring}{测量}\emph{白炽\omterm{def:g3:first-electric-circuit:bulb}{灯泡}}的伏安特性，数据点会偏离
直线：灯丝随电流增大而发热，
炽热的金属\omterm{def:g8:ohms-law:resistance}{电阻}更大——$R$ 并不保持不变。
二极管、\omterm{ex:g6:circuits-and-safety:motor}{电动机}和活体皮肤的特性曲线也会弯曲，
各有其原因。过原点的直线是
\omterm{def:g8:ohms-law:resistor}{电阻器}的特征，而非物质的特征——在用定律之前，
务必先查看元件的特性曲线。
\end{remark}

\section{练习}

\begin{exercise}[$\star$]\label{exo:g8:ohms-law:1}
\omterm{def:g8:ohms-law:resistance}{电阻}\omterm{def:g6:measurement-in-science:measuring}{测量}的是什么？写出它的\omterm{def:g6:measurement-in-science:measuring}{单位}和符号，
并写出\omterm{def:g8:ohms-law:resistance}{欧姆}定律及其三种变形。
\end{exercise}

\begin{exercise}[$\star$]\label{exo:g8:ohms-law:2}
对一个 $30\,\Omega$ 的\omterm{def:g8:ohms-law:resistor}{电阻器}，分别完成以下两组计算：$U =
\qty{6}{V}$，求 $I = {?}$；\ $I = \qty{0.5}{A}$，求 $U = {?}$。
\end{exercise}

\begin{exercise}[$\star$]\label{exo:g8:ohms-law:3}
某元件两端\omterm{def:g8:voltage-voltmeter:voltage}{电压} $U = \qty{4.5}{V}$，此时 $I = \qty{0.09}{A}$。
求它的\omterm{def:g8:ohms-law:resistance}{电阻}。
\end{exercise}

\begin{exercise}[$\star$]\label{exo:g8:ohms-law:4}
在伏安特性曲线实验中，真正\omterm{def:g8:ohms-law:resistor}{电阻器}的各点呈现出什么形状，
较陡的直线又代表什么？
\end{exercise}

\begin{exercise}[$\star$]\label{exo:g8:ohms-law:5}
由本章的表格，验证 $R$ 的数值：依据第二行和
第四行。两次得到相同的数值，
为什么正是整个实验的要点？
\end{exercise}

\begin{exercise}[$\star$]\label{exo:g8:ohms-law:6}
电水壶的加热线圈\omterm{def:g8:ohms-law:resistance}{电阻}为 $25\,\Omega$；另一只电水壶的为
$35\,\Omega$。接在同一市电上，哪一只通过的电流更大，
二者各通过多大的电流（保留一位小数）？
\end{exercise}

\begin{exercise}[$\star$]\label{exo:g8:ohms-law:7}
利用 $U = R I$ 解释：两个\omterm{def:g4:series-circuits:series}{串联}\omterm{def:g3:first-electric-circuit:bulb}{灯泡}共用同一电流，
\omterm{def:g8:ohms-law:resistance}{电阻}较大的\omterm{def:g3:first-electric-circuit:bulb}{灯泡}为何分得较大的\omterm{def:g8:voltage-voltmeter:voltage}{电压}？
\end{exercise}

\begin{exercise}[$\star$]\label{exo:g8:ohms-law:8}
用\omterm{def:g8:ohms-law:resistance}{电阻}的语言描述：\omterm{def:g7:short-circuits-safety:device}{短路}时的电流暴增；
\ 断开\omterm{ex:g3:first-electric-circuit:switch}{开关}时的总阻断。
\end{exercise}

\begin{exercise}[$\star\star$]\label{exo:g8:ohms-law:9}
一只精密指示灯允许通过 \qty{20}{mA} 的电流。
若将它接在 \qty{9}{V} 的\omterm{def:g3:first-electric-circuit:battery}{电池}上，
\omterm{def:g3:first-electric-circuit:bulb}{灯泡}本身约分得 \qty{2}{V} 的\omterm{def:g8:voltage-voltmeter:voltage}{电压}，那么应\omterm{def:g4:series-circuits:series}{串联}多大的\omterm{def:g8:ohms-law:resistor}{电阻器}来保护它？
（先求\omterm{def:g8:ohms-law:resistor}{电阻器}分得的\omterm{def:g8:voltage-voltmeter:voltage}{电压}，再求 $R$。）
\end{exercise}

\begin{exercise}[$\star\star$]\label{exo:g8:ohms-law:10}
实测某白炽\omterm{def:g3:first-electric-circuit:bulb}{灯泡}的伏安特性：（\qty{1}{V}，
\qty{0.25}{A}），（\qty{4}{V}，\qty{0.5}{A}）。证明不存在
单一的 $R$ 能同时符合这两个点，
并说明曲线弯曲背后的物理过程。
\end{exercise}

\begin{exercise}[$\star\star$]\label{exo:g8:ohms-law:11}
两个\omterm{def:g8:ohms-law:resistor}{电阻器} $10\,\Omega$ 与 $20\,\Omega$ \omterm{def:g4:series-circuits:series}{串联}接在
\qty{6}{V} 电源上。利用\omterm{def:g4:series-circuits:series}{串联}规律及各\omterm{def:g8:ohms-law:resistance}{电阻}上的\omterm{def:g8:ohms-law:resistance}{欧姆}定律：
求回路电流及各\omterm{def:g8:ohms-law:resistance}{电阻}两端的\omterm{def:g8:voltage-voltmeter:voltage}{电压}。（提示：
电流 $I$ 满足 $6 = 10 I + 20 I$。）
\end{exercise}

\begin{exercise}[$\star\star\star$]\label{exo:g8:ohms-law:12}
同样两个\omterm{def:g8:ohms-law:resistor}{电阻器}，现改为\omterm{def:g5:series-parallel-circuits:parallel}{并联}接在 \qty{6}{V} 电源上。求各\omterm{def:g5:series-parallel-circuits:parallel}{支路}的电流，
\omterm{def:g3:first-electric-circuit:battery}{电池}的总电流，然后
——漂亮的收尾——求单个\omterm{def:g8:ohms-law:resistance}{电阻} $R$，使其从
\qty{6}{V} 电源获得相同的总电流。将它与两个\omterm{def:g8:ohms-law:resistor}{电阻器}分别比较，
并用道路的比喻说明，为什么两个\omterm{def:g5:series-parallel-circuits:parallel}{并联}的\omterm{def:g8:ohms-law:resistance}{电阻}
合起来的\omterm{def:g8:ohms-law:resistance}{电阻}比其中任何一个都\emph{小}。
\end{exercise}

\section{问题：质量控制检测台}

\begin{problem}[{周末问题——电阻器工厂的检验日；
人像、公差，以及
一个根本不是电阻器的元件}]\label{pb:g8:ohms-law:1}
工厂的质检台检测当天的产品——
还有一个冒牌货。由你来操作仪表。

\medskip\noindent\textbf{第一部分——校准检测台。}
\begin{enumerate}
  \item 画出（或描述）测试电路：电源、
        \omterm{def:g8:intensity-ammeter:ammeter}{电流表}、被测元件、\omterm{def:g8:voltage-voltmeter:voltmeter}{电压表}——谁与之\omterm{def:g4:series-circuits:series}{串联}，
        谁与之\omterm{def:g5:series-parallel-circuits:parallel}{并联}？
  \item 样品一在 \qty{4.5}{V} \omterm{def:g8:voltage-voltmeter:voltage}{电压}下通过 \qty{0.15}{A} 的电流。
        它的\omterm{def:g8:ohms-law:resistance}{电阻}是多少？
  \item 标签标注为 $33\,\Omega$，公差为百分之五。计算可接受的范围，
        并判断
        样品一是否合格。（先求 $33$ 的百分之五。）
  \item 样品二标注为 $150\,\Omega$：预测它在
        \qty{4.5}{V} \omterm{def:g8:voltage-voltmeter:voltage}{电压}下的电流，以便检测台心中有数，
        （保留一位小数，\omterm{def:g6:measurement-in-science:measuring}{单位}用 \unit{mA}）。
\end{enumerate}

\medskip\noindent\textbf{第二部分——完整的人像。}
\begin{enumerate}[resume]
  \item 样品三的数据表：$(1.5, 0.05)$、$(3.0, 0.10)$、
        $(4.5, 0.15)$、$(6.0, 0.20)$——\omterm{def:g6:measurement-in-science:measuring}{单位}为\omterm{def:g8:voltage-voltmeter:voltage}{伏特}、\omterm{def:g8:intensity-ammeter:intensity}{安培}。拟合
        $R$，并指出这些数字中允许取
        单一值的那一特征。
  \item 样品四的数据表：$(1.5, 0.30)$、$(3.0, 0.45)$、
        $(4.5, 0.54)$、$(6.0, 0.60)$。指出这个冒牌货：
        不存在单一的 $R$，并指出偏移的方向。
  \item 样品四实际上是一只小型灯丝\omterm{def:g3:first-electric-circuit:bulb}{灯泡}，
        来自下一条生产线。讲述它弯曲的过程——
        以及为什么检测台的直线测试恰好是一个
        \omterm{def:g8:ohms-law:resistor}{电阻器}检测器。
  \item 检测台的规程写道：“每张人像都必须在
        稳定\omterm{def:g3:temperature-thermometers:temperature}{温度}下拍摄。”这一规定所保护的，是\omterm{def:g8:ohms-law:resistance}{欧姆}
        定律的哪一条款？
\end{enumerate}

\medskip\noindent\textbf{第三部分——应用
部门。}
\begin{enumerate}[resume]
  \item 一份订单要求一只\omterm{def:g8:ohms-law:resistor}{电阻器}将电流限制为
        \qty{30}{mA}，电源为 \qty{12}{V}（负载
        自身分得的\omterm{def:g8:voltage-voltmeter:voltage}{电压}可忽略）。计算应报出的阻值。
  \item 一位客户抱怨其 $25\,\Omega$ 的加热丝线圈
        “只”通过 \qty{9.2}{A} 的电流，而\omterm{def:g7:short-circuits-safety:fuse}{保险丝}允许的是
        $10$。计算其抱怨意味着怎样的电源\omterm{def:g8:voltage-voltmeter:voltage}{电压}，
        并用算术向其解释。
  \item 原型实验室要求工厂提供最粗、最短的铜跳线，
        “\omterm{def:g8:ohms-law:resistance}{电阻}尽可能接近零\omterm{def:g8:ohms-law:resistance}{欧姆}”。
        这种跳线生来是为扮演哪种电路角色？
        若被误用，其接近零的\omterm{def:g8:ohms-law:resistance}{电阻}会使它靠近哪一种
        日常的“大敌”？
  \item 填写检验记录，写下三句话——
        定律、它的人像，以及元件得以签字的
        唯一条件。
\end{enumerate}
\end{problem}
